\documentclass[11pt,letterpaper]{report}

% ============================================================
% Packages
% ============================================================
\usepackage[margin=1in]{geometry}
\usepackage{amsmath, amssymb, amsthm, mathtools}
\usepackage{bm}
\usepackage{graphicx}
\usepackage{float}
\usepackage{booktabs}
\usepackage{longtable}
\usepackage{array}
\usepackage{multirow}
\usepackage{caption}
\usepackage{subcaption}
\usepackage{xcolor}
\usepackage{hyperref}
\usepackage{natbib}
\usepackage{enumitem}
\usepackage{setspace}
\usepackage{fancyhdr}
\usepackage{algorithm}
\usepackage{algpseudocode}

% Hyperref styling
\hypersetup{
  colorlinks=true,
  linkcolor=blue!50!black,
  citecolor=blue!50!black,
  urlcolor=blue!70!black,
}

% Bibliography style
\bibliographystyle{apalike}

% Theorem environments
\newtheorem{definition}{Definition}[chapter]
\newtheorem{proposition}{Proposition}[chapter]
\newtheorem{theorem}{Theorem}[chapter]
\newtheorem{lemma}{Lemma}[chapter]
\newtheorem{remark}{Remark}[chapter]
\newtheorem{assumption}{Assumption}[chapter]

% Custom notation commands (single source of truth)
\newcommand{\R}{\mathbb{R}}
\newcommand{\E}{\mathbb{E}}
\newcommand{\Var}{\operatorname{Var}}
\newcommand{\Cov}{\operatorname{Cov}}
\newcommand{\Corr}{\operatorname{Corr}}
\newcommand{\diag}{\operatorname{diag}}
\newcommand{\tr}{\operatorname{tr}}
\newcommand{\vectr}[1]{\boldsymbol{#1}}
\newcommand{\mat}[1]{\boldsymbol{#1}}
\newcommand{\T}{^{\top}}

% Header/footer
\pagestyle{fancy}
\fancyhf{}
\fancyhead[L]{\leftmark}
\fancyhead[R]{\thepage}
\renewcommand{\headrulewidth}{0.4pt}

\onehalfspacing

% ============================================================
% Title
% ============================================================
\title{
  \vspace{-1cm}
  \textbf{An Industry-Grade Long/Short Fundamental Factor Investing Pipeline for US Equities} \\[0.5em]
  \large Mathematical Framework, Implementation, and Empirical Results
}
\author{
  Amos Anderson \\
  \small Stony Brook University \\
  \small AMS520 --- Machine Learning in Finance
}
\date{Spring 2026}

% ============================================================
% Document
% ============================================================
\begin{document}

\maketitle

\begin{abstract}
\noindent
This report documents the end-to-end construction of a long/short fundamental
factor investing pipeline for US equities, from raw CRSP and Compustat data
through paper trading deployment on Alpaca. We build on the cross-sectional
characteristics of \citet{chen2024deep}, extending their dataset through 2024,
and organize the pipeline in eleven stages: universe construction, factor
engineering, factor validation, signal combination, risk modeling, transaction
cost modeling, portfolio construction, backtesting, attribution, deployment,
and monitoring. A central methodological contribution is the portfolio risk
model, which extends the BARRA-style cross-sectional factor decomposition
$\Sigma = BFB\T + D$ with ARMA-GARCH dynamics on factor returns and Normal
Inverse Gaussian innovations on idiosyncratic residuals. Every stage is
documented with a formal mathematical problem statement, derivations,
references to primary sources, and empirical validation. The pipeline is
implemented as a modular Python package with per-stage notebooks, unit tests,
and configuration artifacts.
\end{abstract}

\tableofcontents
\clearpage

% ============================================================
% Main content
% ============================================================
\chapter{Introduction}
\label{ch:introduction}

\textit{This section introduces the project's motivation, scope, and contribution.
To be written as implementation proceeds.}

\section{Motivation}
\section{Scope and Contribution}
\section{Report Structure}
\section{Reproducibility}
\chapter{Data Foundation}
\label{ch:data_foundation}

\textit{Stage 0 — summary of the CPZ reconstruction, extension to 2024,
and validation against Ken French. To be written after the final audit of
the reconstruction.}

\section{The Chen--Pelger--Zhu Dataset}
\section{Reconstruction from CRSP and Compustat}
\section{Extension Through 2024}
\section{Validation Against CPZ}
\section{External Validation Against Ken French}
\section{The January 2021 Episode and Its Implications}
\section{Summary}
\input{sections/01_universe.tex}
\input{sections/02_factors.tex}
\input{sections/03_validation.tex}
\input{sections/04_alpha.tex}
\input{sections/05_risk.tex}
\input{sections/06_costs.tex}
\input{sections/07_portfolio.tex}
\input{sections/08_backtest.tex}
\input{sections/09_attribution.tex}
\input{sections/10_deployment.tex}
\input{sections/11_monitoring.tex}
\input{sections/12_discussion.tex}

% ============================================================
% Bibliography
% ============================================================
\bibliography{references}

% ============================================================
% Appendices
% ============================================================
\appendix
\chapter{Notation}
\label{app:notation}

Table~\ref{tab:notation} summarizes notation used throughout the report.

\begin{table}[H]
\centering
\caption{Notation reference.}
\label{tab:notation}
\small
\begin{tabular}{ll}
\toprule
Symbol & Meaning \\
\midrule
$i$ & Stock index, $i = 1, \ldots, N_t$ \\
$t$ & Time index (monthly unless noted) \\
$k$ & Factor index, $k = 1, \ldots, K$ \\
$r_{i,t}$ & Excess return of stock $i$ in period $t$ \\
$z_{k,i,t}$ & Value of factor $k$ for stock $i$ at time $t$ (standardized) \\
$B_{i,k,t}$ & Exposure of stock $i$ to factor $k$ at time $t$ \\
$f_{k,t}$ & Return to factor $k$ in period $t$ \\
$e_{i,t}$ & Idiosyncratic return of stock $i$ \\
$\alpha_{i,t}$ & Expected excess return forecast for stock $i$ \\
$w_{i,t}$ & Portfolio weight on stock $i$ at time $t$ \\
$\vectr{w}_t$ & Portfolio weight vector at time $t$ \\
$\mat{B}_t$ & Exposure matrix (stocks $\times$ factors) \\
$\mat{F}_t$ & Factor covariance matrix \\
$\mat{D}_t$ & Idiosyncratic variance matrix (diagonal) \\
$\mat{\Sigma}_t$ & Portfolio covariance: $\mat{\Sigma}_t = \mat{B}_t \mat{F}_t \mat{B}_t\T + \mat{D}_t$ \\
$U_t$ & Investable universe at time $t$ \\
$\text{IC}_t$ & Cross-sectional information coefficient at $t$ \\
$\lambda$ & Risk aversion parameter \\
$G$ & Gross leverage \\
TC & Transaction cost function \\
\bottomrule
\end{tabular}
\end{table}
\chapter{Proofs and Derivations Deferred from Main Text}
\label{app:proofs}

\textit{Proofs and derivations moved from the main text for readability.
Populated as derivations are produced.}

\end{document}